% RESUME_TEMPLATE.tex
% One-page, ATS-friendly resume template (A4/Letter configurable)
% Works for Chinese/English with ctex.
%
% Density preset: slightly denser than v6 to increase info density while staying readable.
% If still >1 page, Skill's self-check loop will trim content first, then apply further density tweaks.

\documentclass[10pt]{ctexart}

\usepackage{geometry}
% NOTE: switch paper size by replacing a4paper -> letterpaper when needed.
% Keep margins modest; do not go below ~1.0cm unless last-resort.
\geometry{a4paper, top=1.1cm, bottom=1.1cm, left=1.3cm, right=1.3cm}

\setlength{\parindent}{0pt}
\setlength{\parskip}{0pt}
\linespread{0.98}

\usepackage{enumitem}
% Compact lists (but avoid 0 spacing that hurts readability).
\setlist[itemize]{leftmargin=*,topsep=1pt,itemsep=0.2pt,parsep=0pt,partopsep=0pt}

\usepackage[hidelinks]{hyperref}
\usepackage{tabularx}
\usepackage{titlesec}
\usepackage{microtype}

% Section title: smaller vertical gaps.
\titleformat{\section}{\normalsize\bfseries}{}{0pt}{}[\vspace{-3pt}\titlerule]
\titlespacing*{\section}{0pt}{4pt}{2pt}

% Header spacing: tighter.
\newcommand{\name}[1]{\begin{center}\LARGE\bfseries #1\end{center}\vspace{-4mm}}
\newcommand{\contact}[1]{\begin{center}\footnotesize #1\end{center}\vspace{-4mm}}

% Entries
\newcommand{\entry}[4]{\textbf{#1} \hfill #2\\ \textit{#3} \hfill \textit{#4}\\}
\newcommand{\proj}[2]{\textbf{#1} \hfill \textit{#2}\\}

\begin{document}

% ====== Header ======
\name{YOUR NAME}
\contact{City \quad|\quad +1-xxx \quad|\quad \href{mailto:you@email.com}{you@email.com}}

% ====== Sections (replace titles per language) ======
\section{Education}
\entry{University}{YYYY--YYYY}{Degree, Major}{City}

\section{Experience}
\entry{Company}{YYYY.MM--YYYY.MM}{Title}{City}
\begin{itemize}
  \item Action + method + result (metric).
\end{itemize}

\section{Projects}
\proj{Project}{YYYY.MM--YYYY.MM}
\begin{itemize}
  \item Built ... (metric).
\end{itemize}

\section{Skills}
\begin{tabularx}{\linewidth}{@{}l X@{}}
  \textbf{Languages} & ... \\
  \textbf{Frameworks} & ... \\
  \textbf{Tools} & ... \\
\end{tabularx}

\end{document}
